\par\medskip\noindent
2.2. Consideremos ahora, además de $P$ y $Q$, un grupo conmutativo
$G$. En (1.1.3) definimos un homomorfismo canónico, funtorial en
$P,Q,G$,
\begin{equation}\tag{2.2.1}
 \Biext^1(P,Q;G)
 \longrightarrow
 \Hom\bigl(\Tor_1(P,Q),G\bigr).
\end{equation}
Por otra parte, si $n$ es un entero positivo (respectivamente, si
$\ell$ es un número primo), el homomorfismo $\alpha^n_{P,Q}$ de
(2.1.1) (respectivamente, el homomorfismo $\alpha^{(\ell)}_{P,Q}$ de
(2.1.13.2)) da un homomorfismo
\[
 \Hom\bigl(\Tor_1(P,Q),G\bigr)
 \longrightarrow
 \Hom({}_nP\otimes{}_nQ,G)
\]
(respectivamente, un homomorfismo
\[
 \Hom\bigl(\Tor_1(P,Q),G\bigr)
 \longrightarrow
 \Hom\bigl(T_\ell(P)\otimes T_\ell(Q),T_\ell(G)\bigr)).
\]
Estos últimos homomorfismos son también funtoriales en $P,Q,G$.
Componiéndolos con (2.2.1), obtenemos homomorfismos canónicos,
funtoriales en $P,Q,G$,
\begin{equation}\tag{2.2.2}
 \varphi^n=\varphi^n_{P,Q}:
 \Biext^1(P,Q;G)
 \longrightarrow
 \Hom({}_nP\otimes{}_nQ,G),
\end{equation}
\begin{equation}\tag{2.2.3}
 \varphi^{(\ell)}=\varphi^{(\ell)}_{P,Q}:
 \Biext^1(P,Q;G)
 \longrightarrow
 \Hom\bigl(T_\ell(P)\otimes T_\ell(Q),T_\ell(G)\bigr).
\end{equation}
Para
\begin{equation}\tag{2.2.4}
 \xi\in\Biext^1(P,Q;G),
\end{equation}

%%%%%%%%%%  scan idx 249  |  folio 238  %%%%%%%%%%

denotemos su imagen por (2.2.2) (respectivamente, por (2.2.3)) con
$\varphi^n_\xi$ (respectivamente, con $\varphi^{(\ell)}_\xi$):
\begin{equation}\tag{2.2.5}
 \varphi^n_\xi:{}_nP\otimes{}_nQ
 \longrightarrow G
 \qquad\text{(o, equivalentemente, }{}_nG\text{)},
\end{equation}
\begin{equation}\tag{2.2.6}
 \varphi^{(\ell)}_\xi:
 T_\ell(P)\otimes T_\ell(Q)
 \longrightarrow T_\ell(G).
\end{equation}

2.2.7. Evidentemente, conocer $\varphi^{(\ell)}_\xi$ equivale a
conocer los $\varphi^n_\xi$ para todos los $n=\ell^\nu$, con
$\nu\geq0$. Para $n$ fijo, conocer $\varphi^n_\xi$ equivale a conocer,
sobre la imagen de
\[
 \sisi{\varphi^n_{P,Q}:{}_nP\otimes{}_nQ
 \longrightarrow\uTor_1(P,Q;G)}%
 {\alpha^n_{P,Q}:{}_nP\otimes{}_nQ
 \longrightarrow\uTor_1(P,Q)},
\]
el homomorfismo
\begin{equation}\tag{2.2.8}
 \uTor_1(P,Q)\longrightarrow G
\end{equation}
que es la imagen de $\xi$ por (2.2.1).

Obsérvese que si $P=nP$ y $Q=nQ$, entonces la imagen precedente es
$\,{}_n\uTor_1(P,Q)$. Esto se deduce fácilmente considerando la
sucesión exacta (1.3.5) y la sucesión exacta análoga para $Q$.
Además, si $P$ y $Q$ son divisibles, es decir, si $nP=P$ y $nQ=Q$
para todo entero positivo $n$, se comprueba fácilmente que
$\uTor_1(P,Q)$ es de torsión, es decir,
\[
 \uTor_1(P,Q)=\varinjlim_n{}_{n}\uTor_1(P,Q).
\]
Así, en este caso, conocer los $\varphi^n_\xi$ para todo $n$ equivale
a conocer (2.2.8), mientras que conocer $\varphi^{(\ell)}_\xi$ para
un primo dado $\ell$ equivale a conocer la restricción de (2.2.8) a
la componente $\ell$-primaria de $\uTor_1(P,Q)$.

Explicitemos ahora las propiedades de simetría de los acoplamientos
\sisi{(2.2.4)}{(2.2.5)} y (2.2.6) asociados a las biextensiones.
Observemos primero que se tiene un diagrama conmutativo
\begin{equation}\tag{2.2.9}
\begin{tikzcd}[column sep=large,row sep=huge]
 \Biext^1(P,Q;G)
   \arrow[r,"\sim"]
   \arrow[d]&
 \Biext^1(Q,P;G)
   \arrow[d]\\
 \Hom\bigl(\uTor_1(P,Q),G\bigr)
   \arrow[r,"\sim"]&
 \Hom\bigl(\uTor_1(Q,P),G\bigr),
\end{tikzcd}
\end{equation}

%%%%%%%%%%  scan idx 250  |  folio 239  %%%%%%%%%%

donde las flechas horizontales son las flechas de simetría (la primera
de ellas se definió en VII 2.7) y las flechas verticales son (2.2.1).
La conmutatividad resulta de la de los dos cuadrados que componen el
diagrama más completo
\begin{equation}\tag{2.2.10}
\begin{tikzcd}[column sep=large,row sep=large]
 \Biext^1(P,Q;G)
   \arrow[r,"\sim"]
   \arrow[d]&
 \Biext^1(Q,P;G)
   \arrow[d]\\
 \Ext^1\!\left(P\mathbin{\overset{\mathbb L}{\otimes}}Q,G\right)
   \arrow[r,"\sim"]
   \arrow[d]&
 \Ext^1\!\left(Q\mathbin{\overset{\mathbb L}{\otimes}}P,G\right)
   \arrow[d]\\
 \Hom\bigl(\uTor_1(P,Q),G\bigr)
   \arrow[r,"\sim"]&
 \Hom\bigl(\uTor_1(Q,P),G\bigr).
\end{tikzcd}
\end{equation}
La primera flecha vertical de cada columna es el isomorfismo canónico
de VII 3.6.5, y la segunda es el homomorfismo evidente. El segundo
cuadrado de (2.2.10) es trivialmente conmutativo: por definición, el
isomorfismo de simetría
\[
 P\mathbin{\overset{\mathbb L}{\otimes}}Q
 \xrightarrow{\ \sim\ }
 Q\mathbin{\overset{\mathbb L}{\otimes}}P
\]
induce el isomorfismo de simetría
$\uTor_1(P,Q)\simeq\uTor_1(Q,P)$ y, para un complejo variable
$L_\bullet$, el isomorfismo canónico
\[
 \Ext^1(L_\bullet,G)
 \simeq
 \Hom\bigl(H_1(L_\bullet),G\bigr)
\]
es funtorial en $L_\bullet$.

En cuanto a la conmutatividad del primer cuadrado de (2.2.10), se
deduce explicitando simultáneamente los isomorfismos de simetría y los
isomorfismos verticales definidos en VII 3.6.5. Dejamos al lector los
detalles de la verificación. Se utiliza el hecho de que el
isomorfismo de simetría
\[
 L_\bullet(P)\otimes L_\bullet(Q)
 \simeq
 L_\bullet(Q)\otimes L_\bullet(P)
\]
induce en grado $1$ el isomorfismo
\[
 L_0(P)\otimes L_1(Q)+L_1(P)\otimes L_0(Q)
 \simeq
 L_0(Q)\otimes L_1(P)+L_1(Q)\otimes L_0(P),
\]

%%%%%%%%%%  scan idx 251  |  folio 240  %%%%%%%%%%

que, sobre los sumandos correspondientes, viene dado por el
isomorfismo de simetría ordinario de los módulos subyacentes no
graduados (y por tanto no introduce ningún signo).

En vista de la conmutatividad de (2.2.9), la anticonmutatividad de
(2.1.10.1), expresada en términos de los $\varphi^n$ o de los
$\varphi^{(\ell)}$, se convierte en la relación siguiente.

\medskip
\noindent\underline{Proposición 2.2.11.}
\underline{Los diagramas siguientes son anticonmutativos:}
\begin{equation}\tag{2.2.11.1}
\begin{tikzcd}[row sep=large,column sep=huge]
 \Biext^1(P,Q;G)
   \arrow[rr,"\sim"]
   \arrow[dd,"{\varphi^n_{P,Q}}"']&&
 \Biext^1(Q,P;G)
   \arrow[dd,"{\varphi^n_{Q,P}}"]\\
 &- &\\
 \Hom({}_nP\otimes{}_nQ,G)
   \arrow[rr,"\sim"]&&
 \Hom({}_nQ\otimes{}_nP,G),
\end{tikzcd}
\end{equation}
\begin{equation}\tag{2.2.11.2}
\begin{tikzcd}[row sep=large,column sep=huge]
 \Biext^1(P,Q;G)
   \arrow[rr,"\sim"]
   \arrow[dd,"{\varphi^{(\ell)}_{P,Q}}"']&&
 \Biext^1(Q,P;G)
   \arrow[dd,"{\varphi^{(\ell)}_{Q,P}}"]\\
 &- &\\
 \Hom\bigl(T_\ell(P)\otimes T_\ell(Q),T_\ell(G)\sisi{}{)}
   \arrow[rr,"\sim"]&&
 \Hom\bigl(T_\ell(Q)\otimes T_\ell(P),T_\ell(G)\bigr),
\end{tikzcd}
\end{equation}
\underline{donde las flechas horizontales son las flechas de simetría.}

Consignemos el corolario importante.

\medskip
\begingroup
\sloppy
\noindent\underline{Corolario 2.2.12.}
\uline{Supongamos que $P=Q$, y que $\xi$ es la clase de una
biextensión simétrica (respectivamente, antisimétrica) de $(P,P)$ por
$G$; es decir, supongamos que ${}^s\xi=\xi$ (respectivamente,
${}^s\xi=-\xi$), donde ${}^s\xi$ denota la imagen de $\xi$ por
simetría. Entonces las aplicaciones bilineales $\varphi^n_\xi$, sobre
${}_nP\times{}_nP$, y $\varphi^{(\ell)}_\xi$, sobre
$T_\ell(P)\times T_\ell(P)$, son antisimétricas (respectivamente,
simétricas).}
\par
\endgroup

%%%%%%%%%%  scan idx 252  |  folio 241  %%%%%%%%%%

\medskip
\noindent\underline{Observación 2.2.13.}
Supongamos que $\xi$ es simétrica. Si $n$ (respectivamente, $\ell$) es
impar, 2.2.12 muestra que $\varphi^n_\xi$ (respectivamente,
$\varphi^{(\ell)}_\xi$) es en realidad una forma alternada. Esta
conclusión no es válida en general sin la restricción de que $n$
(respectivamente, $\ell$) sea impar, como se ve ya cuando $P$ es el
grupo ordinario $\mathbb Z/2\mathbb Z$ y $G=\mathbb Z/2\mathbb Z$.

Supongamos, sin embargo, que $2P=P$. Este es el caso, en particular,
cuando $P$ es un esquema abeliano, trabajando como de costumbre con la
topología fppf o fpqc. Las formas anteriores siguen siendo entonces
alternadas. Basta probar que $\varphi^{2^\nu}_\xi(x,x)=0$. Como
$2P=P$, un elemento dado $x\in{}_{2^\nu}P(S)$ puede escribirse
$x=2y$, con $y\in{}_{2^{\nu+1}}P(S)$. Por consiguiente,
\[
 \varphi^{2^\nu}_\xi(x,x)
 =2\varphi^{2^{\nu+1}}_\xi(y,y)=0,
\]
porque $\varphi^{2^{\nu+1}}_\xi$ es antisimétrica.
